\section{Тестирование}

Тесты реализованы на Google Test в файле \texttt{tests/test\_variant.cpp} и
подключены к CTest через \texttt{gtest\_discover\_tests}. Всего написано
\textbf{26 тестов} при требовании не менее~10, распределённых по трём наборам.

\begin{table}[htbp]
\centering
\caption{Наборы тестов}
\begin{tabular}{@{}l r >{\raggedright\arraybackslash}p{8cm}@{}}
\hline
\textbf{Набор} & \textbf{Тестов} & \textbf{Что проверяет} \\
\hline
\texttt{ArrayOps} & 14 & Создание, освобождение, вставка, удаление, изменение размера \\
\texttt{Variant2} & 6 & Сортировка выбором, разворот, полный алгоритм варианта \\
\texttt{BinarySearch} & 6 & Бинарный поиск, в том числе на дубликатах и пустом массиве \\
\hline
\end{tabular}
\end{table}

\subsection{Покрытые сценарии}

\begin{itemize}
  \item \textbf{Память:} создание массива с нулевой инициализацией, создание
        размера~0 с возвратом \texttt{nullptr}, обнуление указателя после
        \texttt{array\_delete}.
  \item \textbf{Вставка:} в начало, в середину, в конец, за пределами конца
        (с фактической позицией \texttt{size}), в пустой массив.
  \item \textbf{Удаление:} из середины, последнего элемента, позиции вне
        диапазона --- массив и длина не изменяются.
  \item \textbf{Изменение размера:} рост с заполнением новых элементов нулями,
        усечение, переход к размеру~0 с возвратом \texttt{nullptr}, изменение
        размера из \texttt{nullptr}.
  \item \textbf{Алгоритм варианта:} сортировка выбором, сортировка массива
        длины 1 и \texttt{nullptr}, разворот массива чётной длины и длины~1,
        полный алгоритм на контрольном примере, обработка пустого массива.
  \item \textbf{Поиск:} найденное значение в середине и в конце, повторяющиеся
        элементы (проверяется факт нахождения одного из вхождений), отсутствие
        значения ниже, выше и между границами диапазона, пустой массив.
\end{itemize}

\subsection{Контрольный пример}

Алгоритм варианта проверен на примере из задания. Сортировка выбором по
возрастанию, затем разворот:

\begin{center}
\begin{tabular}{|c|c|c|}
\hline
Исходный массив & После сортировки & После разворота \\
\hline
\texttt{\{3, 1, 4, 1, 5\}} & \texttt{\{1, 1, 3, 4, 5\}} & \texttt{\{5, 4, 3, 1, 1\}} \\
\hline
\end{tabular}
\end{center}

\subsection{Результат}

Все 26 тестов из трёх наборов проходят успешно. Сборка выполняется без
предупреждений компилятора при включённых флагах
\texttt{-Wall -Wextra -Wpedantic} (GCC) и \texttt{/W4} (MSVC).
